% Vietnamese translation for OI Public Library
% Source-PDF: IOI中国国家候选队论文/2006/朱泽园/朱泽园.pdf
\documentclass[11pt]{article}
\usepackage{oipl-vn}

\title{Thuật toán mới cho Half-plane Intersection và giá trị thực tiễn}
\author{Zeyuan Zhu}
\date{29 tháng 12 năm 2005}

\begin{document}
\maketitle

\origtitle{New algorithm for Half-plane Intersection and its Practical Value}
\sourcepath{IOI Candidate Team Papers/2006/Zeyuan Zhu/Zeyuan Zhu.pdf}

\begin{center}
  \small Luận văn cho kỳ tuyển chọn đội tuyển Trung Quốc năm 2006

  \smallskip
  Zeyuan Zhu, lớp 12, Nanjing Foreign Language School, Jiangsu, China

  \smallskip
  Tặng người mẹ yêu quý của tôi, Xiaoli Xu.
\end{center}

\noindent\textbf{E-mail:} \texttt{zhuzeyuan[at]hotmail[dot]com}

\noindent\textbf{Từ khóa:} Half-plane, Intersection, Feasible Region, Algorithm,
Polygon, Practical.

\begin{abstract}
Bài viết trình bày một thuật toán \(O(n\log n)\) mới cho
\term{Half-plane Intersection} (viết tắt là HPI), một trong các bài toán được
thảo luận nhiều trong khoa học máy tính. Trọng tâm là giá trị của thuật toán
trong ứng dụng thực tế; ở một mức độ nhất định có thể hạ độ phức tạp xuống
\(O(n)\), trong khi thuật toán vẫn rất dễ cài đặt.

Mục 1 giới thiệu bài toán HPI. Mục 2 chuẩn bị một thuật toán tuyến tính cho
\term{Convex Polygon Intersection} (CPI). Dựa trên đó, mục 3 nhắc lại ngắn gọn
lời giải HPI phổ biến bằng \term{Divide-and-Conquer} (D\&C). Mục 4 trình bày
chi tiết thuật toán mới \term{Sort-and-Incremental} (S\&I). Mục 5 vừa tổng kết
vừa bàn về ứng dụng thực tế và so sánh với thuật toán cũ ở mục 3.
\end{abstract}

\noindent\textbf{Mốc thời gian.} Ý tưởng xuất hiện vào tháng 4 năm 2005; một
phần được cài đặt vào tháng 6 năm 2005\footnote{Bài toán con của HPI xuất hiện
trong kỳ thi USA Invitational Computing Olympiad.}; một bài toán được ra vào
tháng 7 năm 2005\footnote{Tác giả đã ra một bài HPI trên Peking University
Online Judge, kèm giới thiệu ngắn về thuật toán.}; công bố dưới dạng bài viết
trên USENET ngày 6 tháng 11 năm 2005\footnote{\url{http://groups.google.com/group/sci.math/msg/68e9d9fc634f0d92}.}

\section{Giới thiệu}

Một đường thẳng trong mặt phẳng thường được biểu diễn bởi
\[
  ax+by=c.
\]
Tương tự, dạng bất đẳng thức
\[
  ax+by\le c
\]
biểu diễn một nửa mặt phẳng, tức một phía của đường thẳng đó. Cần chú ý rằng
\(ax+by\le c\) và \(-ax-by\le -c\) là hai nửa mặt phẳng đối nhau, khác với
trường hợp đường thẳng \(ax+by=c\) và \(-ax-by=-c\) biểu diễn cùng một đường.

Bài toán \term{Half-plane Intersection} (HPI) được phát biểu như sau: cho
\(n\) nửa mặt phẳng
\[
  a_i x+b_i y\le c_i,\qquad 1\le i\le n,
\]
hãy xác định tập tất cả các điểm thỏa mãn đồng thời \(n\) bất đẳng thức đó.

Miền khả thi, tức giao của các nửa mặt phẳng, có dạng một miền lồi và có thể
không bị chặn. Trong thực hành ta thường thêm bốn nửa mặt phẳng tạo thành một
hình chữ nhật đủ lớn để bảo đảm diện tích sau giao là hữu hạn. Trong các mục
sau, ta giả sử miền khả thi bị chặn và có diện tích hữu hạn.

Mỗi nửa mặt phẳng tạo ra nhiều nhất một cạnh của đa giác lồi kết quả, nên miền
lồi thu được có nhiều nhất \(n\) cạnh. Tuy vậy, giao cũng có thể suy biến
thành một đường thẳng, một tia, một đoạn thẳng, một điểm, hoặc là rỗng.

\section{Convex Polygon Intersection (CPI)}

Bài toán giao hai đa giác lồi \(A\) và \(B\) thành một đa giác lồi mới có thể
giải bằng Line Segment Intersection trong thời gian \(O(n\log n)\), khi tổng
số cạnh là \(O(n)\). Ở đây ta phác thảo một cách đơn giản hơn và hiệu quả hơn:
phương pháp quét mặt phẳng.

Ý tưởng chính là tính các giao điểm của cạnh làm điểm cắt, rồi chia biên của
\(A\) và \(B\) thành các đoạn nằm ngoài và các đoạn nằm trong đa giác kia. Các
đoạn biên nằm trong được nối với nhau, tạo thành đường bao của đa giác giao.
Trong hình minh họa của bản gốc, các đoạn biên trong được tô đậm và chính là
đường bao cần tìm.

Giả sử có một đường quét thẳng đứng đi từ trái sang phải. Các hoành độ được
đường quét xử lý gọi là \term{sự kiện \(x\)}. Tại mọi thời điểm, đường quét có
nhiều nhất bốn giao điểm với hai đa giác đã cho\footnote{Giả sử không có cạnh
nào song song với đường quét. Nếu có, ta có thể xoay mặt phẳng một góc thích
hợp; nếu không xoay, cần xử lý nhiều trường hợp biên.}:
\begin{itemize}
  \item giao với vỏ trên của \(A\), ký hiệu \(A_u\);
  \item giao với vỏ dưới của \(A\), ký hiệu \(A_l\);
  \item giao với vỏ trên của \(B\), ký hiệu \(B_u\);
  \item giao với vỏ dưới của \(B\), ký hiệu \(B_l\).
\end{itemize}

Trên đường quét hiện tại, điểm thấp hơn trong hai điểm \(A_u,B_u\) và điểm cao
hơn trong hai điểm \(A_l,B_l\) tạo thành đoạn của miền giao hiện tại. Nếu đoạn
này không rỗng thì nó nằm trong cả hai đa giác.

Dĩ nhiên ta không thể quét qua mọi hoành độ hữu tỉ. Gọi các cạnh chứa
\(A_u,A_l,B_u,B_l\) lần lượt là \(e_1,e_2,e_3,e_4\). Sự kiện \(x\) kế tiếp
được chọn trong các đầu mút của bốn cạnh này và bốn giao điểm tiềm năng
\[
  e_1\cap e_3,\quad e_1\cap e_4,\quad e_2\cap e_3,\quad e_2\cap e_4.
\]
Thao tác trên có thể cài đặt trong thời gian \(O(n)\), vì chỉ có \(O(n)\) sự
kiện \(x\), còn việc duy trì \(A_u,A_l,B_u,B_l\) tốn \(O(1)\) cho mỗi sự kiện.

\section{Lời giải phổ biến: Divide-and-Conquer (D\&C)}

Cách tiếp cận cơ bản dựa trên chia để trị.
\begin{description}[leftmargin=7em,style=nextline]
  \item[Chia:] Chia \(n\) nửa mặt phẳng thành hai tập có kích thước
  \(\lfloor n/2\rfloor\) và \(\lceil n/2\rceil\).
  \item[Trị:] Tính đệ quy miền khả thi, tức giao, của từng tập con.
  \item[Gộp:] Tính giao của hai đa giác lồi thu được bằng thuật toán CPI tuyến
  tính ở mục 2.
\end{description}

Độ phức tạp thời gian tổng cộng được suy ra từ phương trình truy hồi
\[
  T(n)=2T(n/2)+O(n),
\]
nên
\[
  T(n)=O(n\log n).
\]

\section{Thuật toán mới: Sort-and-Incremental (S\&I)}

Trước hết ta định nghĩa góc cực của một nửa mặt phẳng. Chẳng hạn:
\begin{itemize}
  \item với nửa mặt phẳng dạng \(x-y\ge \text{hằng số}\), góc cực được định
  nghĩa là \(\pi/4\);
  \item với nửa mặt phẳng dạng \(x+y\le \text{hằng số}\), góc cực được định
  nghĩa là \(3\pi/4\);
  \item với nửa mặt phẳng dạng \(x+y\ge \text{hằng số}\), góc cực được định
  nghĩa là \(-\pi/4\);
  \item với nửa mặt phẳng dạng \(x-y\le \text{hằng số}\), góc cực được định
  nghĩa là \(-3\pi/4\).
\end{itemize}
Nói cách khác, ta gắn cho mỗi nửa mặt phẳng một hướng của đường biên để các
nửa mặt phẳng có thể được sắp theo góc cực. Với nửa mặt phẳng dạng
\(ax+by\le c\), ta gọi \(c\) là hằng số của nó.

Thuật toán Sort-and-Incremental được mô tả chi tiết như sau.

\paragraph{Bước 1.}
Chia các nửa mặt phẳng thành hai tập. Một tập có góc cực trong
\((-\pi/2,\pi/2]\); tập còn lại có góc cực trong
\[
  (-\pi,-\pi/2]\cup(\pi/2,\pi].
\]
Thực ra chỉ cần chia thành hai khoảng độ dài \(\pi\), chẳng hạn
\((a,a+\pi]\) và \((a+\pi,a+2\pi]\). Cách chia trên được dùng để giải thích rõ
hơn.

\paragraph{Bước 2.}
Xét tập có góc cực trong \((-\pi/2,\pi/2]\); tập còn lại được xử lý tương tự ở
bước 3 và bước 4. Sắp xếp các nửa mặt phẳng theo góc cực. Với các nửa mặt
phẳng có cùng góc cực, chỉ giữ lại nửa mặt phẳng chặt nhất, tức nửa mặt phẳng
bị các nửa mặt phẳng còn lại phủ lên; các nửa mặt phẳng khác có thể xóa đi.
Nếu tất cả được viết cùng dạng \(ax+by\le c\) với cùng hướng, đó là nửa mặt
phẳng có hằng số nhỏ nhất.

\paragraph{Bước 3.}
Bắt đầu với hai nửa mặt phẳng có góc cực nhỏ nhất sau khi sắp xếp. Tính giao
điểm hai đường biên của chúng và đẩy cả hai nửa mặt phẳng vào một ngăn xếp.

Sau đó, mỗi lần thêm một nửa mặt phẳng mới theo thứ tự góc cực tăng dần. Tính
giao điểm giữa nửa mặt phẳng mới và nửa mặt phẳng ở đỉnh ngăn xếp. Nếu giao
điểm này nằm bên phải giao điểm của hai nửa mặt phẳng trên cùng trong ngăn
xếp, ta loại đỉnh ngăn xếp. Tiếp tục kiểm tra với hai nửa mặt phẳng mới ở
đỉnh ngăn xếp; nếu giao điểm hiện tại vẫn nằm bên phải giao điểm ở đỉnh, lại
loại tiếp. Lặp cho đến khi giao điểm hiện tại nằm bên trái giao điểm ở đỉnh,
rồi đẩy nửa mặt phẳng hiện tại vào ngăn xếp.

Hình minh họa trong bản gốc cho thấy trực giác của thao tác này: khi thêm một
nửa mặt phẳng có góc cực lớn hơn, các giao điểm cũ nằm bên trái giao điểm mới
không còn tạo thành phần biên hợp lệ của nửa đa giác đang xây dựng, nên các
nửa mặt phẳng tương ứng bị pop khỏi ngăn xếp.

\paragraph{Bước 4.}
Các giao điểm của những cặp nửa mặt phẳng kề nhau trong ngăn xếp tạo thành
một nửa của đa giác lồi. Với hai tập ở bước 1, ta thu được hai nửa của đa giác:
\((-\pi/2,\pi/2]\) cho vỏ trên, còn
\((-\pi,-\pi/2]\cup(\pi/2,\pi]\) cho vỏ dưới. Có thể gộp chúng trong thời gian
tuyến tính bằng CPI ở mục 2, nhưng ta còn có thể tận dụng cách lặp sau.

Dựa trên cấu trúc hàng đợi hai đầu, ta liên tục loại một đầu mút đã biết là
không khả thi. Ban đầu bốn con trỏ \(p_1,p_2,p_3,p_4\) trỏ tới các cạnh trái
nhất và phải nhất của vỏ trên và vỏ dưới. Hai con trỏ \(p_1,p_3\) đi sang
phải, còn \(p_2,p_4\) đi sang trái. Tại mọi thời điểm, nếu giao điểm trái nhất
không thỏa nửa mặt phẳng do \(p_1\) hoặc \(p_3\) cung cấp, giao điểm trái nhất
đó chắc chắn phải bị xóa; con trỏ tương ứng đi sang cạnh kề bên phải. Tương
tự, ta xử lý giao điểm phải nhất với \(p_2\) và \(p_4\).

Mỗi lần chỉ thử xóa một giao điểm trái nhất hoặc phải nhất. Lý do là có thể
giao điểm ngoài cùng vi phạm một nửa mặt phẳng, trong khi giao điểm kế bên lại
nằm đúng phía của một nửa mặt phẳng khác và không thể xóa cùng lúc. Lặp quá
trình này cho đến khi không còn cập nhật nào. Trong quá trình đó, số giao điểm
của một vỏ riêng lẻ có thể giảm xuống \(0\), nhưng miền khả thi tổng thể vẫn
có thể không rỗng.

Ngoại trừ bước sắp xếp ở bước 2, mọi phần của thuật toán S\&I đều tuyến tính,
tức \(O(n)\). Nếu dùng quick-sort cho bước 2, độ phức tạp tổng cộng là
\(O(n\log n)\), với hằng số ẩn khá nhỏ.

\section{Giá trị thực tiễn và cách tiếp cận tuyến tính}

Những ý tưởng lớn cần cả cánh để bay lẫn bộ phận hạ cánh để đáp xuống thực tế.
Thuật toán S\&I có vẻ cùng độ phức tạp thời gian với D\&C, nhưng khi cài đặt
nó có một số ưu thế rõ rệt.

\begin{itemize}
  \item Chương trình S\&I dễ viết hơn nhiều so với chương trình D\&C. Bản cài
  đặt bằng C++ dài chưa tới 3 KB.

  \item Hằng số ẩn trong độ phức tạp của S\&I nhỏ hơn đáng kể so với D\&C, vì
  ta không còn cần \(O(n\log n)\) phép tính giao điểm. Nói chung, chương trình
  S\&I chạy nhanh hơn chương trình D\&C xấp xỉ năm lần.

  \item Phép tính giao điểm là thao tác rất quan trọng trong cả hai thuật toán
  vì nó chậm, tương đương hàng trăm phép cộng. CPI, thuật toán chuẩn bị được
  D\&C gọi nhiều lần, cần \(O(n)\) phép tính giao điểm ở mỗi lần gộp, khiến
  tổng thời gian tăng lên. Ngược lại, S\&I chỉ tính giao điểm \(O(n)\) lần
  trong mọi trường hợp.

  \item Nếu các nửa mặt phẳng đầu vào đều nằm trong \((-\pi/2,\pi/2]\), hoặc
  trong bất kỳ khoảng góc nào có độ dài \(\pi\), chương trình S\&I có thể được
  rút ngắn đáng kể, chỉ khoảng hai mươi dòng C++. D\&C không có lợi thế tương
  tự. Một bài toán trong kỳ USA Invitational Computing Olympiad đã khai thác
  đúng tình huống này.
\end{itemize}

\paragraph{Tối ưu tiệm cận xuống tuyến tính.}
Nút thắt của thuật toán là sắp xếp. Cần chú ý rằng phép sắp xếp ở đây không
nhất thiết là sắp xếp so sánh. Khóa của các phần tử cần sắp là góc cực, có thể
xác định bởi hệ số góc, tức một phân số thập phân. Do đó ta có thể thay
quick-sort \(O(n\log n)\) bằng radix-sort \(O(n)\), và hạ độ phức tạp tiệm cận
của thuật toán xuống \(O(n)\). Dẫu vậy, trong thực tế cách \(O(n)\) thường có
thể chạy chậm hơn cách \(O(n\log n)\) vì cần thêm bộ nhớ.

\paragraph{Cảm nghĩ về sáng tạo.}
Một phát minh không thuộc về người chỉ nảy ra ý tưởng. Phần lớn mọi người đều
có ý tưởng. Khác biệt giữa một người bình thường và một nhà phát minh là vì lý
do nào đó nhà phát minh tin vào chính mình và có thôi thúc đưa ý tưởng đến kết
quả. Henri Matisse, họa sĩ và nhà điêu khắc người Pháp, từng nói rằng không
có gì khó hơn đối với một họa sĩ thật sự sáng tạo hơn là vẽ một bông hồng, bởi
trước khi làm được điều đó, người ấy phải quên tất cả những bông hồng từng
được vẽ. Mang theo cả tinh thần đổi mới lý tưởng lẫn thực tiễn, tôi đang trên
đường. Còn bạn thì sao?

\begin{thebibliography}{9}
\bibitem{mehlhorn}
Kurt Mehlhorn.
\newblock \textit{Data Structures and Algorithms, Vol. 3}.
\newblock Springer-Verlag, New York, 1984.

\bibitem{clrs}
Thomas H. Cormen, Charles E. Leiserson, Ronald L. Rivest, Clifford Stein.
\newblock \textit{Introduction to Algorithms}, 2nd edition.
\newblock MIT Press, New York, 2001.

\bibitem{barequet}
Gill Barequet.
\newblock Computational Geometry, Lecture 4, Spring 2004/2005.

\bibitem{telikepalli}
Kavitha Telikepalli.
\newblock Algorithms and Data Structures, Lecture 3, Winter 2003/2004.
\end{thebibliography}

\end{document}
